\chapter{加与减：正交宇宙物理速览}

本章从一条极简单的公式出发：只需把时空几何定律中的一个减号换成一个加号，
便能得到一个物理定律迥然不同的世界——\emph{黎曼宇宙}（Riemannian universe）。
我们将速览这一改动所引发的连锁后果：从固有的时间与光速，到能量与动量，
再到光的静止质量与物质的稳定性。读者会看到一些似曾相识的现象、一些奇异而优美
的结构，也会看到一些潜藏的危险。

\section{从减号到加号：时空中的勾股定理}

一百多年来我们已经知道，理解我们宇宙中时间的最佳方式，是把时间与熟悉的空间
三维结合起来，构成一个服从自身独特几何定律的四维\emph{时空}（space-time）。
牛顿（Newton）把时间视为绝对而普遍的量：它以完全相同的速率、整齐划一地为每个人
向前推进；而自爱因斯坦（Einstein）以来我们认识到，时间流逝所度量的，是我们在
四维时空中所走路径的一个属性，它与空间中一条路径的\emph{长度}非常相似。

没有人会对这一点感到惊讶：空间中连接同样两点的不同路径可以有不同的长度。
如果我从珀斯（Perth）走最直接的路线开车到悉尼（Sydney），而你中途绕道达尔文
（Darwin），那么我们的里程表读数不同，这毫不奇怪。我们现在明白，时空中的
“绕道”对时间流逝也有类似的作用。例如，若你留在地球上，“旅行”于 2050 年的
除夕与 2060 年的除夕之间；而我虽然同样在两次除夕与你相聚，却在中间这段时间
往返了一次半人马座$\alpha$（Alpha Centauri），那么我将老去约六年，而你老去十年。

空间中的绕道使路径更长，时空中的绕道却使它更短。其原因可以表述得很简单：
在度量时间间隔时，勾股定理的“时空版”用的是\emph{差}而不是\emph{和}——
是减号而不是加号。

勾股定理指出，直角三角形斜边的平方等于另外两条直角边平方之和。例如：

\orthofig{images/00__PM.001.png}{空间中的勾股定理：斜边平方等于两直角边平方之和}

而时空版则把平方之和换成平方之差：

\orthofig{images/00__PM.002.png}{勾股定理的时空版：用平方之差取代平方之和}

为使公式简洁，我们选取一种让时间与空间处于同等地位的单位：时间用“年”，
距离用“光年”。可以把时间间隔的平方写成

\begin{equation}
(\Delta \tau)^2 = (\Delta t)^2 - (\Delta x)^2 , \tag{1}
\end{equation}

其中 $\Delta\tau$ 为固有时（proper time）。利用两个这样的三角形，即可验证
前面关于半人马座$\alpha$之行的说法——它约有四光年之遥，而我此行的结果可能是
我只老去六年，而留在家中的你老去十年：

\orthofig{images/00__PM.003.png}{往返半人马座$\alpha$的时空三角形}

单程 $\Delta x = 4$ 光年、$\Delta t = 5$ 年，于是

\[
(\Delta \tau)^2 = 5^2 - 4^2 = 9 \;\Longrightarrow\; \Delta\tau = 3\ \text{年},
\]

往返两程合计 $6$ 年；而留在地球的人经历 $10$ 年。

当然，这些三角形在我们看来“比例不对”：那条本应长“$3$ 年”的斜边，
被画成明显长于代表“$5$ 年”的竖直边。但这正是要点所在：时空的“洛伦兹几何”
（Lorentzian geometry）并不等同于空间的几何，因此不可能在一张纸上把每个细节
都忠实地画出来。

现在可以很轻松地道出我们自己的宇宙与我一直在探索的替代宇宙之间的区别了。
在黎曼宇宙中，普通的勾股定理\emph{总是}成立，即使涉及时间间隔时也是如此。
这个宇宙的时空图可以完全忠实地画在纸上！那种朴素的直觉——认为一条“$5$ 年”
的边与一条“$4$ 光年”的边构成的三角形应当画成\emph{这样}：

\orthofig{images/00__PM.004.png}{黎曼宇宙中时空三角形的勾股关系}

如今是完全正确的。在黎曼宇宙中，绕道去往一颗遥远的恒星，对旅行者而言
总是意味着经历的时间更多，而不是更少。

星际航行对旅行者比对其他人耗时更长，听起来或许只是恼人；但由这同一处
简单改动所派生的差异远不止于此，它们波及从夜空的面貌到物质的稳定性等方方面面。

在我们宇宙中，(1) 式里的减号意味着：一旦距离（以光年计）超过时间（以年计），
平方之差就会变成一个负数。例如，设想用 $4$ 年穿行 $5$ 光年，公式给出的历时
平方便是 $-9$。既然没有实数的平方为负，这样一次假想的行程显然在性质上是不同的。
而这实际上是在告诉我们：没有任何东西能在如此短的时间内跨越如此大的距离。

在黎曼宇宙中，与之对应的行程毫无问题。对 $\Delta x = 5$ 光年、家乡系中
$\Delta t = 4$ 年的航行，船员测得的历时应为

\[
\Delta\tau = \sqrt{4^2 + 5^2} = \sqrt{41} \approx 6.4\ \text{年},
\]

而家乡的人只经历 $4$ 年，飞船被看作以光速的一又四分之一倍行进：

\[
v = \frac{5}{4}\,c = 1.25\,c .
\]

在黎曼宇宙中，不存在普适的速度上限。

事实上，黎曼宇宙中也不存在普适的\emph{光}速，因此“光年”和“光速”这些词
实际上只在我们自己的宇宙中才恰当。在黎曼宇宙中，代替“光速”的，是一个把时间
换算为距离的换算因子，它使勾股公式呈现出简洁漂亮的形式。正如你若用“公里”量
南北方向的距离、用“英里”量东西方向的距离，勾股定理就会失效；在黎曼宇宙中，
同样需要通过实验去测定：时间与距离各取什么单位，才能使该定理成立。

\section{没有普适光速的宇宙}

生活在一个光以不同速度传播的宇宙中会是什么样子？当然，光穿过玻璃之类的
透明介质时，在我们宇宙中也会这样；但在黎曼宇宙中，即使在真空里也会如此。
我们把不同波长的光感知为不同的颜色；在黎曼宇宙中，这些不同的波长将\emph{总是}
以不同的速度传播。

要找出波长与速度的关系，可以注意这样一点：无论光速如何，光波前在\emph{时空}
中的间距总是相同的。为什么？我们期望自己宇宙中的物理不会偏向三维空间里的任何
方向：若我们用一支激光笔朝东南方发射，我们会认为它产生的光，与朝正北或西北方
发射时产生的光性质完全相同。同样，黎曼宇宙中底层的物理也应当对刻画光速的
那个时空方向无动于衷。

假设一束速度适中的光脉冲的时空图如下：

\orthofig{images/00__PM.005.png}{时空中的光脉冲与波前}

这里我们为光脉冲画出的“路径”显示了它随时间的整体空间运动。正如水波的波峰
垂直于波传播的方向，标记这束光振荡峰值的波前，也垂直于它的路径。

光的\emph{波长}是这些波前在空间中的间距。当然，我们称为“空间”的那个方向
取决于我们自身的运动状态，因此我们测得的波长也会随之改变。相反，波前沿其
法向（即与波前垂直的方向）度量的\emph{间距}，是光本身的性质，与谁在测量无关。

要得到较慢光脉冲的波长，只需把前一幅图中的整族波前整体转动，
而不改变它们彼此之间的间距：

\orthofig{images/00__PM.006.png}{时空中较慢的光脉冲与波前}

由此可见，传播得越慢的光，波长越长。把它翻译成我们感知各种波长颜色的语言，
便可以说：红光将比紫光慢。

当你仰望黎曼宇宙的夜空时，从每颗恒星看到的红光，要比紫光多花数年乃至数百
年才能抵达你这里——所以除非这颗恒星完全没有横向运动，不同颜色将来自天空上
稍有差异的位置。于是每颗恒星都不再表现为单一的光点，其像会被展宽成一段
光谱，或一条“色谱拖尾”（colour trail），其中红、紫两端之间的张角直接度量了
该恒星的横向速度。

\begin{notebox}[补充：波长与速度的关系]
若记波前族的法向间距为 $d$（这是不变量），并设光在 $(x,t)$ 平面内的传播方向为
$(v,1)$，则在黎曼（欧氏）几何中，波前法向的单位分量为
$n_x = v/\sqrt{1+v^2}$。沿空间 $x$ 轴量得的波长为
\[
\lambda = \frac{d}{n_x} = \frac{d\sqrt{1+v^2}}{v},
\]
可见速度 $v$ 越小，波长 $\lambda$ 越长——与上面“转动波前族”的几何论证完全一致。
（此推导为补充说明，原文只给出转动波前族的定性论证。）
\end{notebox}

\section{能量-动量矢量}

除了没有普适的速度上限，关于黎曼宇宙中物质的行为还有什么可说的？要理解最
引人注目的差异，我们得先勾勒我们自己宇宙物理学中的几个简单概念。

对于运动的物体，有两个有用的量：它的动量（momentum）与动能（kinetic energy）。
刻画它们的最简单方式，也许是设想把一颗子弹射入某种屏障，该屏障以恒定力阻碍
其运动使之停下。子弹的\emph{动量}决定它停下来的快慢，而它的\emph{动能}决定它
停下来之前走过多远。这两个量都只取决于子弹的质量与速度。

自相对论问世以来，我们认识到：正如空间与时间最好被看成同一实体（时空）的
两个方面，运动物体的能量与动量也是同一个几何对象的两个方面，这个对象叫作
\emph{能量-动量矢量}（energy-momentum vector）。这听起来或许很专业，但实际上
它非常容易想象。对于静止的物体，能量-动量矢量就是一支长度等于物体质量、
指向时空\emph{时间}方向的箭头。

\orthofig{images/00__PM.007.png}{我们宇宙中的能量-动量矢量}

若使同一物体运动起来，原来的箭头只是被倾斜，以反映物体在给定时间内穿越的
空间大小：慢速物体略微倾斜，快速物体倾斜更大。

在上图中，箭头看上去似乎不只是被倾斜，还变长了。但这只是把时空画在纸面上
造成的错觉；若我们用时空版的勾股定理去计算倾斜箭头的长度，会发现它与那支竖直
箭头长度完全相同。

这些能量-动量矢量有什么用？若你测量物体能量-动量矢量在时间方向上的投影长度，
它告诉你物体的\emph{能量}。若你测量它在空间方向上的投影长度，它告诉你物体的
\emph{动量}。

\orthofig{images/00__PM.008.png}{我们宇宙中的总能量、动量与动能}

这里唯一稍显棘手之处在于：能量-动量矢量在时间方向上的投影给出的是物体的
\emph{总能量}（total energy），它不仅包含动能，还包含与物体质量相联系的能量。
因此，静止物体虽然有零动量（正如你所预期），却\emph{不是}零总能量；它的总能量
完全来自其质量（如前所述，质量决定了这些箭头的长度）。当同一物体被赋予运动时，
其总能量在“静质量能”之上的\emph{增量}，就是它的动能。

为什么我们没有引用那个著名的静质量能公式 $E=mc^2$？因为我们所用的这类单位里，
光速 $c$ 总等于 $1$（例如一年每光年），于是 $E=m$，无需 $c^2$。

那么在黎曼宇宙中又如何？我们仍然画一支长度等于物体质量的箭头，并根据物体的
速度给它一个相应的倾斜：

\orthofig{images/00__PM.009.png}{黎曼宇宙中的能量-动量矢量}

这一次，倾斜看起来就是普通的倾斜：我们肉眼就能看出两支箭头长度相同。时空的
特殊勾股定理已被普通勾股定理取代，因此我们可以在普通的纸上画出黎曼宇宙的
能量-动量矢量，并用普通的直尺量其长度。

这一改动如何影响运动物体的能量与动量？

\orthofig{images/00__PM.010.png}{黎曼宇宙中的总能量、动量与动能}

在黎曼宇宙中，倾斜箭头使它在时间方向上的投影\emph{小于}物体静止时的投影。
所以物体运动时的总能量比静止时\emph{更小}！我们仍然可以把物体的动能定义为
二者之差，但它现在代表的是一种能量的减少，而非增加。

声称运动物体比静止物体能量更低，听起来极其古怪；但事实上，动能与总能量之间
这种“上下颠倒”的特征，在最简单的情形里根本不会显现。我们熟悉的全部经典牛顿
力学，对慢速运动的物体仍然适用：你完全可以在黎曼宇宙里打一局台球，而察觉不到
任何不同。低速物体的动量与动能的近似公式，在两个宇宙中完全相同——只是动能
差一个负号，而台球桌上发生的任何事都不会揭示这个负号的意义。动能相对静质量能
是负的，而在我们宇宙中静质量能只在核反应中才显现。所以，黎曼宇宙这个古怪的
特性也许会在日常世界中一直隐藏着，只在玄奥的高能实验中才露出端倪？

\begin{equation}
E = \begin{cases}
\sqrt{m^2 + p^2} \approx m + \dfrac{p^2}{2m} & \text{(洛伦兹型)},\\[2ex]
\sqrt{m^2 - p^2} \approx m - \dfrac{p^2}{2m} & \text{(黎曼型)} .
\end{cases}
\end{equation}
式中 $p^2/(2m)=\half mv^2$ 即经典动能。两式在低速下完全对应，仅动能项的符号相反。

\section{光有静止质量}

但有一件远非玄奥的事情被我们忽略了：光！在我们自己的宇宙中，光没有静止质量，
而时空版的勾股定理允许光的能量-动量矢量具有相应的长度\emph{零}。为此，总能量
不必为零，只需等于动量（在光速为 $1$ 的单位下）即可。这样，矢量在空间方向与
时间方向上的投影，便在平方之差里相互抵消了。

\orthofig{images/00__PM.011.png}{我们宇宙中光的能量-动量矢量}

\begin{equation}
E^2 - p^2 = 0 \;\Longrightarrow\; E = p . \tag{2}
\end{equation}

在黎曼宇宙中，这是不可能的：一个矢量长度为零的唯一方式，是它在所有方向上的
投影都为零。静止质量为零就意味着完全没有能量与动量。所以在黎曼宇宙中，
光具有静止质量，和别的一切东西一样。

\begin{equation}
E^2 + p^2 = m^2 \;\Longrightarrow\; m>0 .
\end{equation}

\section{发光的能量学与物质的稳定性}

现在，任何产生光的过程都需要从别的形式（如动能或化学能）获取能量。但当动能
相对总能量为\emph{负}时，“获取”一词就不太恰当了。产生光时我们必须遵守能量
守恒定律，而既然光的总能量为正，那么产生光的那个体系的总能量就必须下降……
也就是说，它的动能必须\emph{增加}。

换言之，一个产生光的体系将经历动能的增加：它的各组成部分最终会比之前运动得
更快！

\orthofig{images/00__PM.012.png}{产生光的反应中的能量平衡}

举一个简单而具体的例子。假设我们做一个产生光的化学反应。当把化学物质混合在
一起时，它们大致是静止的（它们因热而带有一些动能，但我们一开始假定这可以
忽略）。为简单起见，我们还假定这两种化学物质质量相等。于是，在“反应前”的图中，
它们的能量-动量矢量只是两支相同的、指向时间方向的箭头。

“反应后”的图则需要包含所生成的光的能量-动量矢量，以及化学反应产物的能量-动量
矢量。我们假定产物与反应物质量相同；这未必成立，但正如我们将会看到的，
它当然是可能的。

要在“反应后”的图中把所有矢量安排得使总能量不变，产物的能量-动量矢量就必须
被倾斜。也就是说，产物必须运动起来：它们需要获得动能并降低总能量，以便为所
生成的光“垫付”总能量。

\begin{notebox}[反应前后能量配平的代数]
设两种反应物质量相等，均为静止（动量为零）：反应前总能量 $E_{\text{前}} = 2m$。
反应后设两产物质量仍为 $m$、携带反向动量 $\pm p$，并生成总能量为 $E_{\text{光}}>0$
的光。能量守恒要求
\[
E_{\text{前}} = E_{\text{后}} = E_{\text{光}} + 2\sqrt{m^2 - p^2} = 2m .
\]
由 $E_{\text{光}}>0$ 立得 $\sqrt{m^2 - p^2} < m$，即产物必须处于运动状态：其总能量
低于静质量能，从而为光提供了能量。（其中光与产物动量的分配由动量守恒
$\sum \mathbf{p} = \bzero$ 决定。）
\end{notebox}

化学反应中产生的动能表现为热，因此该反应将同时产生光与热。在我们宇宙中，
燃烧任何燃料也会产生光与热，但在那种情况下，它们是以化学能的改变为代价换来的，
而这又表现为产物与反应物之间（极其微小的）质量差异。而在这里，我们无需依赖任何
这样的变化，就制造出了光与热！我们特意令产物的质量与反应物完全相同，但即使把
产物质量取得略高或略低，反应也仍然能够发生。

在黎曼宇宙中，产生光的过程可以产生热，无论是否伴随化学能的改变。当然，化学能
有时是比热更有用的产物。黎曼宇宙中的植物不会吸收光来进行光合作用，反而会
\emph{发出}光，以储存由此生成的化学能。

产生光总是伴随动能的产生，这使黎曼宇宙成为一个危险的地方。生命需要掌握某些
精细的反应，才能安全地从这一过程中获取能量；就连最简单的物质基本单元也将面临
不稳定的威胁。在我们宇宙中，原子或分子可以吸收一个光子而进入一个结合得较松的
高能态；例如阳光中的紫外线可以把某些分子打散。在黎曼宇宙中，同类过程不再需要
外部辐射源：处于最紧束缚态的物质可以\emph{发出}自己的光，并借此获得动能来把
自身撕裂。由于这一效应，只有某些结构是稳定的。

\section{研究异宇宙的价值}

研究一个由不同于我们的定律所统治的宇宙，有什么价值？我们当然可以设想这样
一个宇宙或许真实存在——那里也有居民，也在思索：如果他们自己的某个加号变成了
减号，生活会变成什么样。但把我们的定律拆解开来，看看它们在新假设下如何运作，
也能为我们自己宇宙的物理学提供洞见。通过努力想象事物可能有何不同，我们可以
加深对许多事物的理解——无论是一个人类社会、一颗行星的地质、生命的化学，
还是最基本的物理定律。

黎曼宇宙充满了危险与惊奇。本文只是触及了它的皮毛，但希望读者已从中尝到了
一点它所能带来的乐趣。

\begin{keybox}[本章要点]
\begin{itemize}
  \item 把时空勾股定理的一个减号改为加号：洛伦兹型给出
        $(\Delta\tau)^2=(\Delta t)^2-(\Delta x)^2$，即“绕道使固有时更短”；
        黎曼型给出 $(\Delta\tau)^2=(\Delta t)^2+(\Delta x)^2$，即“绕道使时间更长”。
        时间与距离的换算单位必须由实验确定。
  \item 黎曼宇宙没有普适光速：波长越长（越红）的光传播越慢。波前在时空中的
        \emph{法向间距}是不变量，快慢之别仅表现为整族波前的转动；恒星因而呈现
        “色谱拖尾”，红紫张角直接度量其横向速度。
  \item 能量-动量矢量：长度等于质量，静止时指向时间方向；时间分量是总能量，
        空间分量是动量，动能 = 总能量 $-$ 静质量能。黎曼宇宙中运动物体的总能量
        比静止时\emph{更小}，故动能相对总能量为负；低速极限下与牛顿力学一致，
        仅动能差一个负号。
  \item 黎曼几何中零长矢量只能是零矢量，故光必有静止质量。于是“产生光必须付出
        能量”，发光体系的动能反而增加（表现为放热）；植物的光合作用逆转为
        “发光储能”；处于最紧束缚态的物质可自发光并借此获得动能、撕裂自身，
        使得只有特定结构才稳定。
\end{itemize}
\end{keybox}
